\section{A greedy bound}\label{sec:greedy}

Throughout this section, $G$ is the complete bipartite graph $K_{n,n}$ with vertex classes $A$ and $B$, equipped with a proper $n$-edge-colouring $c\colon E(G)\to[n]$. For a matching $M$ we write $V(M)$ for the set of vertices covered by $M$ and $c(M)=\{c(e):e\in M\}$ for its set of colours.

\begin{definition}\label{def:rainbow}
A matching $M$ in $G$ is \term{rainbow} if $\abs{c(M)}=\abs{M}$. A rainbow matching is \term{maximal} if it is not contained in a larger rainbow matching.
\end{definition}

Since the colouring is proper and uses $n$ colours on the $n^2$ edges of $G$, each colour class is a perfect matching:
\begin{equation}\label{eq:colour-class}
\abs{c^{-1}(x)}=n \quad\text{and}\quad c^{-1}(x)\ \text{is a matching, for every } x\in[n].
\end{equation}

\begin{lemma}\label{lem:extend}
Let $M$ be a rainbow matching in $G$ and let $x\in[n]\setminus c(M)$. If $2\abs{M}<n$, then some edge $e$ of colour $x$ is disjoint from $V(M)$, and $M+e$ is a rainbow matching.
\end{lemma}

\begin{proof}
Suppose every edge of colour $x$ meets $V(M)$. By \eqref{eq:colour-class} the edges of colour $x$ are pairwise disjoint, so each of the $2\abs{M}$ vertices of $V(M)$ lies on at most one of them. Hence $n=\abs{c^{-1}(x)}\leq 2\abs{M}$, a contradiction. So some edge $e$ of colour $x$ avoids $V(M)$; then $M+e$ is a matching, and it is rainbow because $x\notin c(M)$.
\end{proof}

\begin{theorem}\label{thm:greedy}
Every proper $n$-edge-colouring of $K_{n,n}$ contains a rainbow matching with at least $\lceil n/2\rceil$ edges. Equivalently, every Latin square of order $n$ has a partial transversal of size at least $\lceil n/2\rceil$.
\end{theorem}

\begin{proof}
Let $M$ be a maximal rainbow matching. If $\abs{M}=n$ there is nothing to prove. Otherwise some colour $x\in[n]$ does not appear on $M$, and, as $M$ is maximal, Lemma~\ref{lem:extend} gives $2\abs{M}\geq n$. Since $\abs{M}$ is an integer, $\abs{M}\geq\lceil n/2\rceil$.
\end{proof}

\begin{remark}\label{rem:switch}
The proof of Lemma~\ref{lem:extend} only uses one unused colour. If $2\abs{M}$ is close to $n$, then for \emph{every} colour $x\notin c(M)$ almost all edges of colour $x$ have exactly one endpoint in $V(M)$. Koksma's bound~\cite{koksma1969} comes from exploiting this: one swaps an edge of $M$ for an edge of an unused colour and then extends, which is the simplest instance of the switching arguments used in~\cite{hatami2008,kpsy2022}.
\end{remark}

We end the section with the classical example, mentioned in Section~\ref{sec:intro}, showing that a full transversal need not exist.

\begin{claim}\label{claim:cyclic}
If $n$ is even, the cyclic Latin square $L(i,j)=i+j \bmod n$ on $\Z_n$ has no transversal.
\end{claim}

\begin{proof}
Suppose the cells $(i,\sigma(i))$, $i\in\Z_n$, form a transversal, where $\sigma$ is a permutation of $\Z_n$. Then the symbols $i+\sigma(i)$ run over all of $\Z_n$, so, working modulo $n$,
\[
\sum_{s\in\Z_n}s=\sum_{i\in\Z_n}\bigl(i+\sigma(i)\bigr)=2\sum_{i\in\Z_n}i .
\]
Now $\sum_{s=0}^{n-1}s=\frac n2(n-1)\equiv\frac n2 \pmod n$, since $n-1$ is odd, while $2\sum_i i=n(n-1)\equiv0\pmod n$. As $\frac n2\not\equiv0\pmod n$, this is a contradiction.
\end{proof}
